\section{Introduction}
\label{sec:Introduction}

This document gives a brief overview of what \LaTeX\ (pronounced ``Lay-tek") is and highlights some of the advantages it has over Microsoft Word. In a sentence, \LaTeX\ is a well-established piece of typesetting software for producing technical documents where the input is given in plain text which is then compiled into a PDF. I believe that using \LaTeX\ will improve upon several weaknesses in our document production and document control processes and is in line with our digital philosophy of using technology to improve our workflow. It is understandable why one may be wary of any potential barriers to those less comfortable with technology, but the payoff is more than worth it. The main advantages are as follows:

\begin{itemize}
	\item No need to worry about formatting - template is followed automatically.
	\item Equations, cross-references, and citations are far easier than in Word.
	\item Compatible with text-based version control systems such as Git.
	\item The \LaTeX\ compiler knows how to typeset better than you do so documents look professional.
	\item Less buggy than Word.
	\item Vector image compatible and are given by reference.
\end{itemize}

Below I have included the code used to create this section as an example of what \LaTeX\ may look like to a typical user\footnote{Note that this is done by reference to a file so it will update whenever this introduction is changed. This ``single source of truth" functionality is common in \LaTeX.}. Pressing the compile button converts this into a PDF.